The construction of the drought impact database consists of five main stages: data collection, preprocessing, LLM-based information extraction, geospatial processing and aggregation, and validation. Together, these stages transform unstructured newspaper articles into a structured and spatially explicit drought impact database.

This chapter introduces the data collection process, preprocessing steps, and the LLM-based extraction framework. The extraction framework is subsequently evaluated to identify the most suitable model configuration. The extracted impact records are then geocoded, aggregated to NUTS-2 and NUTS-3 regions, and validated using external data sources to assess the quality of the resulting database.

Figure~\ref{fig:pipeline_overview} provides an overview of the complete database construction framework and the corresponding chapters in Part~I.


\begin{figure}[h]
    \centering
    \fbox{
        \parbox{0.96\textwidth}{
        \centering

        \textbf{Input}\\[0.15cm]
        Newspaper Articles

        \vspace{0.3cm}
        $\downarrow$
        \vspace{0.3cm}

        Data Collection
        $\rightarrow$
        Preprocessing
        $\rightarrow$
        Extraction Framework\\
        \small (Chapter~4)

        \vspace{0.3cm}
        $\downarrow$
        \vspace{0.3cm}

        Extraction Performance \& Configuration Selection\\
        \small (Chapter~5)

        \vspace{0.3cm}
        $\downarrow$
        \vspace{0.3cm}

        Geospatial Processing \& Aggregation
        $\rightarrow$
        Visualization \& Validation\\
        \small (Chapter~6)

        \vspace{0.3cm}
        $\downarrow$
        \vspace{0.3cm}

        \textbf{Output}\\[0.15cm]
        Dutch Drought Impact Database
        }
    }
    \caption{Overview of the drought impact database construction framework and the corresponding chapters in Part~I.}
    \label{fig:pipeline_overview}
\end{figure}


% The construction of a drought impact database requires transforming unstructured information sources into structured and spatially explicit impact records. This chapter presents the framework developed to construct the drought impact database used throughout Part~I of this thesis. The workflow consists of five main stages: data collection, preprocessing, LLM-based information extraction, geospatial processing and aggregation, and validation.

% The workflow begins with data collection, where relevant newspaper articles are gathered using a semi-automated approach. These articles are subsequently processed through several preprocessing steps to remove irrelevant information, reduce redundancy, and prepare the text for automated extraction.

% Following preprocessing, large language models (LLMs) are used to extract structured drought impact information from the articles. Chapter~\ref{ch:extraction_performance} investigates how different LLM configurations and structured output constraints influence extraction performance. Selecting an appropriate extraction configuration is challenging because drought impacts described in text can have multiple valid interpretations, requiring evaluation against human judgement.

% After extraction, the identified impact records undergo geospatial processing, where extracted locations are converted into geographic coordinates and aggregated to NUTS-2 and NUTS-3 regions. This step creates a spatially consistent dataset suitable for further analysis and forecasting applications in part II. Finally, the resulting database is analysed through a visualization framework and evaluated using external data sources to assess extraction quality and identify remaining uncertainties.

% The final outcome of this workflow is a structured drought impact database containing spatially explicit impact records, together with an assessment of its quality and limitations.


% \begin{figure}[h]
%     \centering
%     \fbox{
%         \parbox{0.9\textwidth}{
%         \centering
%         Newspaper Articles 
%         $\rightarrow$
%         Data Collection 
%         $\rightarrow$
%         Preprocessing 
%         $\rightarrow$
%         LLM-based Extraction 
%         $\rightarrow$
%         Geospatial Processing \& Aggregation 
%         $\rightarrow$
%         Visualization \& Validation
%         }
%     }
%     \caption{Overview of the drought impact database construction framework. The framework transforms unstructured newspaper articles into structured and spatially explicit drought impact records.}
%     \label{fig:pipeline_overview}
% \end{figure}










% %1. preprecossing
% %2. extraction
% %2. Geocoding
% %4. visuliazation and validaiton

% The 








% \begin{figure}[h]
%     \centering
%     \fbox{
%         \parbox{0.9\textwidth}{
%         \centering
%         Preprocessing $\rightarrow$ LLM extraction $\rightarrow$ Geocoding $\rightarrow$ Visualisation  $\rightarrow$ 
%         Validation
        
%         }
%     }
%     \caption{High-level overview of the drought impact extraction pipeline}
%     \label{fig:pipeline_overview}
% \end{figure}


% To construct and get our drought impact database, several steps have been taken to get to the final results. First we have to collect the right data, this has been been done in the data collection step, wich was done using an semi-automation. Then after that we applied several preporcceing techniche's, to clean the data and make it machine ready. In chapter 5 we further explore llmn configuration settings, and pick an model that works the best for use case, based on human validation. This is not an trivilial task, becauce, one event might have multiple correct anwsers. After Evaluating the best peformance setup, we geoceode the data, and create an aggregation framework, to aggereate our data into nuts3 regoin. WHne that is done we feed the data through an visulization tool, and peform validation studies, using external data sources. This leaves us with an dataset where reliabilty has sort of been estimated towards succes. 







% The proposed system consists of a modular pipeline that transforms unstructured news articles into a structured and geospatially explicit dataset of drought impacts. This system architecture is the core of Part \hyperref[part A]{I}. In the current chapter, we focus on the preprocessing step and the LLM-based extraction framework. While Chapter \ref{ch:extraction_performance} is focused on the Extraction Performance and Configuration Selection. Chapter \ref{ch:geocoding}, we evaluate the final 3-steps of the pipeline Geocoding, visualization and validation.

% This architecture integrates multiple processing stages, including preprocessing, large language model-based information extraction, geocoding, and visualisation, as illustrated in Figure~\ref{fig:pipeline_overview}.

% The pipeline begins with data entry and extraction, followed by preprocessing steps that include cleaning, deduplication, and basic feature extraction (see Section \ref{ch:preprocessing}). The processed text is then passed to a large language model, which identifies and structures drought-related impact information from unstructured articles (see Section \ref{Extraction Framework Design}).

% Extracted location mentions are subsequently transformed into geospatial coordinates through a geocoding step, enabling spatial interpretation of impact events. The resulting structured dataset is exported for downstream analysis (see Section \ref{nuts_coding}).

% A Streamlit-based interface is used for visualisation of geocoded impacts and supports exploratory analysis as well as validation of extracted results (see Section \ref{vis}).
